% Copy-ready front matter for the post-closure Version 2 manuscript.
% Do not include this file in the canonical manuscript until the normalized
% second-moment theorem and the final event assembly have been validated on one
% integrated branch.

\begin{abstract}
Let $G_n\sim G(n,1/2)$.  We prove that the difference between the chromatic
number $\chi(G_n)$ and the cochromatic number $\zeta(G_n)$ is at least a
positive constant times $n/(\log n)^3$ with probability tending to one.  The
argument is uniform across the full rounding phase of the independence-number
threshold.  Its first-moment component compares ordinary colorings with signed
cocoloring witnesses supported on four consecutive class sizes.  For the
second moment, we derive an exact sign-compatibility identity, encode the large
overlap cells by a matching with local deficits, sum the associated physical
matching fibers exactly, and control the residual even-subgraph factor by
injective restriction outside the exposed matching.  A bounded-differences
argument then amplifies the resulting positive-probability seed to a
high-probability cocoloring.
\end{abstract}

\maketitle

\section{Introduction}\label{sec:introduction-v2}

All graphs in this paper are finite, simple, and undirected.  A
\emph{cocoloring} of a graph $G$ is a partition of $V(G)$ into nonempty classes,
each of which is either an independent set or a clique.  The least number of
classes in such a partition is the \emph{cochromatic number} $\zeta(G)$.  This
is the standard invariant introduced by Lesniak and Straight
\citep{lesniak-straight-1977}.  We write $\chi(G)$ for the chromatic number.

Let $G_n\sim G(n,1/2)$ be the labeled random graph on $[n]$ in which the
$\binom n2$ possible edges occur independently with probability $1/2$.  An
event holds \emph{with high probability} if its probability tends to one as
$n\to\infty$.

Erd\H{o}s and Gimbel asked whether
\[
  \chi(G_n)-\zeta(G_n)\longrightarrow\infty
\]
with high probability \citep[p.~263]{erdos-gimbel-1993}.  The problem was
later restated by Gimbel \citep[Section~7.4]{gimbel-2016} and is cataloged as
Erd\H{o}s Problem~625 \citep{bloom-erdos625}.  The main result gives a
full-sequence lower bound of the conjectured polynomial order.

\begin{theorem}\label{thm:main-v2}
For $G_n\sim G(n,1/2)$,
\[
  \mathbb P\!\left(
    \chi(G_n)-\zeta(G_n)
    \ge
    \frac{(\log 2)^2}{8}
    \log\!\left(\frac{1000}{639}\right)
    \frac{n}{(\log n)^3}
  \right)
  \longrightarrow 1.
\]
In particular, $\chi(G_n)-\zeta(G_n)$ tends to infinity with high probability
along the full sequence of integers $n$.
\end{theorem}

The full-sequence assertion is significant.  Near a jump of the
independence-number threshold, changing the relevant class size by one changes
the feasible profile.  A result proved only for a dense set of phase values
therefore need not extend to every integer $n$.  Our estimates are uniform in
the complete phase parameter, including sequences approaching either endpoint
of a phase interval.

\subsection{Background}

The chromatic number of dense random graphs has been studied since the work of
Grimmett and McDiarmid \citep{grimmett-mcdiarmid-1975}.  The first-order
asymptotic was established by Bollob\'as \citep{bollobas-1988}; later
refinements include \citet{mcdiarmid-1990},
\citet{panagiotou-steger-2009}, and \citet{heckel-2018}.  The cochromatic
number fits into the theory of generalized chromatic numbers for hereditary
graph properties developed by Scheinerman and by Bollob\'as--Thomason
\citep{scheinerman-1992,bollobas-thomason-1995}.

For the chromatic--cochromatic difference, Heckel and, independently, Steiner
obtained the first quantitative evidence toward divergence
\citep{heckel-2024-question,steiner-2024}.  Heckel subsequently proved a much
larger lower bound on a phase-dependent set containing approximately $95\%$ of
the integers \citep{heckel-2025-difference}.  The remaining difficulty is to
control the complete phase uniformly.  The present proof uses the signed
first-moment gain and rare-seed amplification from that work, but replaces the
pairwise sign bound by an exact overlap identity and treats every phase with
one four-size profile.

\subsection{Main ideas}

The proof has three stages.

First, we compare two continuous first-moment roots.  The ordinary root governs
proper colorings.  The second root governs a four-size signed witness in which
each class is marked either independent or complete.  The $2^k$ possible marks
create a uniform entropy gain and move the signed root to the left by order
$n/(\log n)^3$.

Second, we prove that this signed root is populated.  For two signed witnesses,
the overlap matrix determines an exact local reward and a binary cycle-space
factor.  Cells whose multiplicity exceeds half the phase cap form a matching.
After fixing this block-level matching, each selected cell is described by its
full endpoint multiplicity and a nonnegative deficit.  The local physical
matching fibers are summed exactly, and the only nonlocal change is one global
falling-factorial ratio.  The remaining even-subgraph contribution is bounded
directly by restricting outside the exposed matching.

Third, Paley--Zygmund gives a positive-probability signed witness.  A
one-Lipschitz induced cocolorable-capacity variable and a simultaneous
leftover-coloring estimate amplify this seed to a high-probability upper bound
for $\zeta(G_n)$.  Intersecting that event with the ordinary-coloring lower
bound completes the proof.

\subsection{Organization of the paper}

Section~2 resolves the full independence-number phase.  Sections~3 and~4
locate the ordinary first-moment root and derive an unrestricted lower bound
for $\chi(G_n)$.  Section~5 constructs the four-size signed profile and proves
the root separation.  Section~6 derives the exact signed-overlap identity, and
Section~7 controls partial diagonals.  Sections~8 and~9 treat the high-cell
matching and the residual attachment, respectively.  Section~10 amplifies the
signed seed, and the final section assembles the two high-probability events.

Exact finite identities, deterministic inequalities, asymptotic estimates, and
probabilistic conclusions are stated separately throughout.  This separation
is important in Sections~8 and~9, where a local product formula and a global
falling-factorial normalization occur in the same calculation.
